\begin{table}[htbp]
\centering
\caption{The contamination on genuinely stuck days}
\label{tab:openingstuck}
\small
\begin{threeparttable}
\begin{tabular}{@{}lrrrr@{}}
\toprule
Opening digit & Stuck days (\%) & Mobile days (\%) & Difference & Stock-days, stuck \\
\midrule
0 & 16.44 & 15.24 & +1.20 & 1,532 \\
1 & 15.47 & 14.39 & +1.07 & 738 \\
2 & 14.08 & 14.05 & +0.03 & 660 \\
3 & 13.24 & 13.98 & -0.74 & 661 \\
4 & 12.80 & 13.87 & -1.07 & 665 \\
5 & 12.78 & 13.80 & -1.02 & 935 \\
6 & 11.89 & 13.87 & -1.97 & 661 \\
7 & 11.69 & 13.73 & -2.04 & 624 \\
8 & 13.78 & 14.15 & -0.37 & 735 \\
9 & 14.98 & 14.15 & +0.83 & 628 \\
\bottomrule
\end{tabular}
\begin{tablenotes}[flushleft]\footnotesize
\item Mean $M^{0}$ by the last digit of the opening price. Stuck means the day's high--low range, scaled by the opening price, falls in the bottom decile of that day's cross-section -- the closest available analogue of Ohta's Table 2 conditioning. The signature of the mechanical contamination is a positive difference at digit zero and negative differences in the middle digits: a stuck day prints mostly its opening neighbourhood, whatever digit that happens to be.
\end{tablenotes}
\end{threeparttable}
\end{table}
